% GENERATED FILE. Edit canonical lessons, then run npm run generate:books.
% canonical-source-hash: fnv1a64:d76964d7812caba5
% canonical-lessons: GU-C318-raja, GU-C318-rani, GU-C318-mahel, GU-C318-darbar, GU-C318-gadi

\chapter{The King and the Queen}
\label{ch:gu-the-king-and-the-queen}
\hlchaptermodality{318}

\begin{chapteropening}
\textbf{By the end of this chapter:} \emph{I can tell who ruled long ago and how they lived.}

Complete the last lesson of chapter 318: I can tell who ruled long ago and how they lived.
\end{chapteropening}

\section[\texorpdfstring{\gu{રાજા} (rājā)}{રાજા (rājā)}]{\gu{રાજા} (rājā) — a king}

\label{lesson:GU-C318-raja}



\begin{quote}
Before the new one: say the Gujarati for the past, then the Gujarati for to be born.
\end{quote}

\subsection*{You'll want to know: \gu{રાજા}}
\textbf{\gu{રાજા}} — \emph{rājā} — \textquotedblleft{}a king\textquotedblright{}.

\begin{quote}
\gu{જૂના} \gu{વખતમાં} \gu{અહીં} \gu{એક} \gu{દયાળુ} \gu{રાજા} \gu{રહેતો} \gu{હતો}.
\end{quote}

\begin{quote}
\emph{jūnā vakhatmā̃ ahī̃ ek dayāḷu rājā raheto hato.}
\end{quote}

\begin{quote}
Long ago a kind king lived here.
\end{quote}

A masculine noun.

\begin{cousinweb}
\textbf{\gu{રાજા}} (\emph{rājā}) is inherited from Sanskrit \emph{rājan}, \textquotedblleft{}a chieftain, a king\textquotedblright{} (CDIAL 10679). It ends in \emph{-ā} but does not change like \textbf{\gu{છોકરો}}: \textbf{\gu{રાજાનું} \gu{ઘર}}, the king's house.
\end{cousinweb}

\begin{grammarlens}[title={What you've built}]
One more word for telling what happened.
\end{grammarlens}

\subsection*{Your turn}
\begin{itemize}
\raggedright
  \item \emph{Say it:} \emph{rājā}
  \item \emph{Say it:} \emph{rājā}, once more
  \item \emph{Say it:} \emph{jūnā vakhatmā̃ ahī̃ ek dayāḷu rājā raheto hato.} — Long ago a kind king lived here.
  \item \emph{Recall:} say the Gujarati for the past, then the Gujarati for to be born, then say \emph{rājā} again
\end{itemize}

\begin{tcolorbox}[breakable,colback=teal!4,colframe=teal!35!black,title={Before you move on}]
What does \gu{રાજા} mean? (\textquotedblleft{}A king\textquotedblright{}.) Say it once more.
\end{tcolorbox}
\section[\texorpdfstring{\gu{રાણી} (rāṇī)}{રાણી (rāṇī)}]{\gu{રાણી} (rāṇī) — a queen}

\label{lesson:GU-C318-rani}



\begin{quote}
Before the new one: say the Gujarati for to be born, then the Gujarati for a king.
\end{quote}

\subsection*{You'll want to know: \gu{રાણી}}
\textbf{\gu{રાણી}} — \emph{rāṇī} — \textquotedblleft{}a queen\textquotedblright{}.

\begin{quote}
\gu{રાણીએ} \gu{ગરીબ} \gu{લોકોને} \gu{અનાજ} \gu{આપ્યું}.
\end{quote}

\begin{quote}
\emph{rāṇīe garīb lokone anāj āpyũ.}
\end{quote}

\begin{quote}
The queen gave grain to the poor.
\end{quote}

A feminine noun.

\begin{cousinweb}
\textbf{\gu{રાણી}} (\emph{rāṇī}) is inherited from Sanskrit \emph{rājñī}, \textquotedblleft{}a queen\textquotedblright{}, and CDIAL lists the Gujarati word (CDIAL 10692). With \textbf{\gu{એ}} she is the doer: \textbf{\gu{રાણીએ} \gu{અનાજ} \gu{આપ્યું}}.
\end{cousinweb}

\begin{grammarlens}[title={What you've built}]
One more word for telling what happened.
\end{grammarlens}

\subsection*{Your turn}
\begin{itemize}
\raggedright
  \item \emph{Say it:} \emph{rāṇī}
  \item \emph{Say it:} \emph{rāṇī}, once more
  \item \emph{Say it:} \emph{rāṇīe garīb lokone anāj āpyũ.} — The queen gave grain to the poor.
  \item \emph{Recall:} say the Gujarati for to be born, then the Gujarati for a king, then say \emph{rāṇī} again
\end{itemize}

\begin{tcolorbox}[breakable,colback=teal!4,colframe=teal!35!black,title={Before you move on}]
What does \gu{રાણી} mean? (\textquotedblleft{}A queen\textquotedblright{}.) Say it once more.
\end{tcolorbox}
\section[\texorpdfstring{\gu{મહેલ} (mahel)}{મહેલ (mahel)}]{\gu{મહેલ} (mahel) — a palace}

\label{lesson:GU-C318-mahel}



\begin{quote}
Before the new one: say the Gujarati for a king, then the Gujarati for a queen.
\end{quote}

\subsection*{You'll want to know: \gu{મહેલ}}
\textbf{\gu{મહેલ}} — \emph{mahel} — \textquotedblleft{}a palace\textquotedblright{}.

\begin{quote}
\gu{રાજા} \gu{અને} \gu{રાણી} \gu{મોટા} \gu{મહેલમાં} \gu{રહેતાં} \gu{હતાં}.
\end{quote}

\begin{quote}
\emph{rājā ane rāṇī moṭā mahelmā̃ rahetā̃ hatā̃.}
\end{quote}

\begin{quote}
The king and queen lived in a big palace.
\end{quote}

A masculine noun.

\begin{cousinweb}
\textbf{\gu{મહેલ}} (\emph{mahel}) is usually traced, through Persian, to Arabic \emph{maḥall}, \textquotedblleft{}a place, a residence\textquotedblright{}. A king and a queen together take the plural \textbf{\gu{રહેતાં} \gu{હતાં}}, the form for a man and a woman.
\end{cousinweb}

\begin{grammarlens}[title={What you've built}]
One more word for telling what happened.
\end{grammarlens}

\subsection*{Your turn}
\begin{itemize}
\raggedright
  \item \emph{Say it:} \emph{mahel}
  \item \emph{Say it:} \emph{mahel}, once more
  \item \emph{Say it:} \emph{rājā ane rāṇī moṭā mahelmā̃ rahetā̃ hatā̃.} — The king and queen lived in a big palace.
  \item \emph{Recall:} say the Gujarati for a king, then the Gujarati for a queen, then say \emph{mahel} again
\end{itemize}

\begin{tcolorbox}[breakable,colback=teal!4,colframe=teal!35!black,title={Before you move on}]
What does \gu{મહેલ} mean? (\textquotedblleft{}A palace\textquotedblright{}.) Say it once more.
\end{tcolorbox}
\section[\texorpdfstring{\gu{દરબાર} (darbār)}{દરબાર (darbār)}]{\gu{દરબાર} (darbār) — a royal court}

\label{lesson:GU-C318-darbar}



\begin{quote}
Before the new one: say the Gujarati for a queen, then the Gujarati for a palace.
\end{quote}

\subsection*{You'll want to know: \gu{દરબાર}}
\textbf{\gu{દરબાર}} — \emph{darbār} — \textquotedblleft{}a royal court\textquotedblright{}.

\begin{quote}
\gu{રાજાના} \gu{દરબારમાં} \gu{એક} \gu{વેપારી} \gu{આવ્યો}.
\end{quote}

\begin{quote}
\emph{rājānā darbārmā̃ ek vepārī āvyo.}
\end{quote}

\begin{quote}
A merchant came to the king's court.
\end{quote}

A masculine noun.

\begin{cousinweb}
\textbf{\gu{દરબાર}} (\emph{darbār}) is usually traced to Persian \emph{darbār}, \textquotedblleft{}a court, a royal audience\textquotedblright{}: the hall where a king hears his people. \textbf{\gu{દરબારમાં}}, at court.
\end{cousinweb}

\begin{grammarlens}[title={What you've built}]
One more word for telling what happened.
\end{grammarlens}

\subsection*{Your turn}
\begin{itemize}
\raggedright
  \item \emph{Say it:} \emph{darbār}
  \item \emph{Say it:} \emph{darbār}, once more
  \item \emph{Say it:} \emph{rājānā darbārmā̃ ek vepārī āvyo.} — A merchant came to the king's court.
  \item \emph{Recall:} say the Gujarati for a queen, then the Gujarati for a palace, then say \emph{darbār} again
\end{itemize}

\begin{tcolorbox}[breakable,colback=teal!4,colframe=teal!35!black,title={Before you move on}]
What does \gu{દરબાર} mean? (\textquotedblleft{}A royal court\textquotedblright{}.) Say it once more.
\end{tcolorbox}
\section[\texorpdfstring{\gu{ગાદી} (gādī)}{ગાદી (gādī)}]{\gu{ગાદી} (gādī) — a throne}

\label{lesson:GU-C318-gadi}



\begin{quote}
Before the new one: say the Gujarati for a palace, then the Gujarati for a royal court.
\end{quote}

\subsection*{You'll want to know: \gu{ગાદી}}
\textbf{\gu{ગાદી}} — \emph{gādī} — \textquotedblleft{}a throne\textquotedblright{}.

\begin{quote}
\gu{રાજાના} \gu{મૃત્યુ} \gu{પછી} \gu{તેનો} \gu{દીકરો} \gu{ગાદી} \gu{પર} \gu{બેઠો}.
\end{quote}

\begin{quote}
\emph{rājānā mṛtyu pachī teno dīkro gādī par beṭho.}
\end{quote}

\begin{quote}
After the king's death his son took the throne.
\end{quote}

A feminine noun.

\begin{cousinweb}
\textbf{\gu{ગાદી}} (\emph{gādī}) comes from \emph{*garda}, a form CDIAL reconstructs, which lists Gujarati \emph{gādī}, \textquotedblleft{}a pad, a raised seat\textquotedblright{} (CDIAL 4053): the cushioned seat of a ruler. \textbf{\gu{ગાદી} \gu{પર} \gu{બેઠો}}, he came to the throne.
\end{cousinweb}

\begin{grammarlens}[title={What you've built}]
One more word for telling what happened.
\end{grammarlens}

\subsection*{Your turn}
\begin{itemize}
\raggedright
  \item \emph{Say it:} \emph{gādī}
  \item \emph{Say it:} \emph{gādī}, once more
  \item \emph{Say it:} \emph{rājānā mṛtyu pachī teno dīkro gādī par beṭho.} — After the king's death his son took the throne.
  \item \emph{Recall:} say the Gujarati for a palace, then the Gujarati for a royal court, then say \emph{gādī} again
\end{itemize}

\begin{tcolorbox}[breakable,colback=teal!4,colframe=teal!35!black,title={Before you move on}]
What does \gu{ગાદી} mean? (\textquotedblleft{}A throne\textquotedblright{}.) Say it once more.
\end{tcolorbox}
